\label{chap:p3-83-46-godel_turing}

\section{Decidibilidad efectiva de cada ventana finita de multiplicidades proyectivas}

\begin{proposition}
Fijada una profundidad \(N\), si los dominios son finitos y los observables
atómicos son computables, toda fórmula del lenguaje finito
\(L_N\) es semánticamente decidible.
\end{proposition}

\begin{proof}
Las fórmulas atómicas se calculan; las conectivas se deciden por tablas de
verdad; cada cuantificador recorre un dominio finito. Se concluye por
inducción sobre la complejidad sintáctica.
\end{proof}

La decidibilidad de cada ventana no proporciona necesariamente un solo
algoritmo que decida una propiedad de la torre completa.

\section{Torre efectiva asociada al cómputo de multiplicidades proyectivas}

Sea \(T_e\) la máquina de Turing de índice \(e\). Para cada \(n\), definimos
\[
 X_n^e
 =
 \{b_n\}
 \cup
 \{c_n:T_e\text{ no se ha detenido en sus primeros }n\text{ pasos}\}.
\]
Los enlaces conservan las etiquetas existentes:
\[
  p_{n+1,n}(b_{n+1})=b_n,
  \qquad
  p_{n+1,n}(c_{n+1})=c_n.
\]
La torre visible posee un solo símbolo \(v_n\), y
\[
  q_n(b_n)=q_n(c_n)=v_n.
\]

\begin{theorem}[Multiplicidad global y parada]
\label{p3-83-c46-s1:thm:halt-fibre}
\[
 \#q_\infty^{-1}(\boldsymbol v)
 =
 \begin{cases}
 1,&T_e\text{ se detiene},\\
 2,&T_e\text{ nunca se detiene}.
 \end{cases}
\]
Por tanto, no existe un algoritmo total que decida la unicidad de la fibra
global para todas las torres efectivas de esta familia.
\end{theorem}

\begin{proof}
La rama \(\boldsymbol b=(b_n)\) existe siempre. Si \(T_e\) se detiene en el
paso \(h\), los estados \(c_n\) desaparecen a partir de \(h\) y no forman una
sección global. Si nunca se detiene, \(\boldsymbol c=(c_n)\) es una segunda
sección. Un decisor de la unicidad resolvería \(\HALT\), contradiciendo
\citep{turing1936}.
\end{proof}

Definimos
\[
  U_e:\#q_\infty^{-1}(\boldsymbol v)=1.
\]
Entonces \(U_e\iff T_e\downarrow\).

\begin{theorem}[Incompletitud uniforme]
\label{p3-83-c46-s1:thm:uniform-incompleteness}
Sea \(\mathsf T\) una teoría recursivamente axiomatizable que formaliza las
torres \(X^e\) y es semánticamente correcta para \(U_e\) y
\(\neg U_e\). Entonces \(\mathsf T\) no decide todos los \(U_e\).
\end{theorem}

\begin{proof}
Si decidiera cada \(U_e\), para un índice \(e\) se enumerarían en paralelo
las pruebas de \(U_e\) y \(\neg U_e\) hasta encontrar una. Por corrección,
la primera aparecería exactamente cuando la máquina se detiene y la segunda
cuando no. Se obtendría un decisor de \(\HALT\).
\end{proof}

\begin{figure}[H]
\centering
\begin{tikzpicture}[>=Latex,node distance=11mm and 15mm,
  box/.style={rounded corners,draw=hmtblue,fill=hmtlightblue,
    minimum width=39mm,minimum height=13mm,align=center},
  warn/.style={rounded corners,draw=hmtred,fill=hmtlightred,
    minimum width=39mm,minimum height=13mm,align=center}]
  \node[box] (sim) {simulación finita de \(T_e\)};
  \node[box,below=of sim] (levels) {\(X_n^e\) computable\\para cada \(n\)};
  \node[box,below=of levels] (fibre)
    {\(\#q_\infty^{-1}(v)=1\iff T_e\downarrow\)};
  \node[warn,below left=of fibre] (undec) {sin decisor uniforme};
  \node[warn,below right=of fibre] (inc)
    {teoría efectiva correcta\\incompleta para los \(U_e\)};
  \draw[->,thick,hmtblue] (sim)--(levels);
  \draw[->,thick,hmtblue] (levels)--(fibre);
  \draw[->,thick,hmtred] (fibre)--(undec);
  \draw[->,thick,hmtred] (fibre)--(inc);
\end{tikzpicture}
\caption{De la decisión finita a la incompletitud uniforme.}
\label{p3-83-c46-s1:fig:godel-turing}
\end{figure}

\section{Completitud finita de los niveles HMT y dominio exacto de la obstrucción Gödel--Turing}

El nombre académico utilizado es:
\[
\boxed{\text{Teorema HMT de incompletitud uniforme para multiplicidades
proyectivas}.}
\]
También se emplea la denominación abreviada «tercer resultado de Gödel en
HMT». La prueba
es Gödel--Turing: no reproduce la diagonal autorreferente de
\cite{goedel1931}, sino que reduce HALT a la multiplicidad de una fibra.

\begin{remark}[Condiciones de validez]
El teorema no prueba inconsistencia de ZFC. Tampoco impide demostrar
directamente que una torre particular posea una sección única. Impide un
procedimiento universal de unicidad para todas las torres efectivas.
\end{remark}

% REMATE_LOCAL_GODEL_CONCLUSION
\Needspace{25\baselineskip}
\section{Relación exacta con las \(9\) fases del Topological Phase Kernel}

Las \(9\) fases del TPK aportan un algoritmo concreto que, con sus lectores declarados,
particiona el bloque de refinamiento en \(729\) subcilindros y selecciona uno de ellos
en cada etapa finita.
El paso a la sección infinita usa un teorema global de no vaciedad,
compatibilidad y unitariedad.

El resultado Gödel--Turing dice:
\[
\text{no existe un decisor de unicidad para toda torre efectiva}.
\]
No dice:
\[
\text{ninguna torre concreta puede tener una prueba uniforme propia}.
\]
Y una prueba particular no proporciona el decisor universal prohibido.


\begin{remark}[Conclusión estructural]
Decisión finita, existencia del límite, unicidad de una instancia y decisión
uniforme de todas las instancias son cuatro niveles. El cierre holográfico
exige declararlos en lugar de trasladar automáticamente una propiedad de un
nivel a otro.
\end{remark}

\begin{remark}[Separación de resultados y alcance]
La torre reductora, la equivalencia \(U_e\iff T_e\downarrow\) y el teorema
metamatemático son \emph{resultados establecidos en las secciones precedentes}. La posición central dentro
del cierre del infinito es una organización expositiva que no altera el contenido lógico ni el dominio del teorema.
\end{remark}
